\scanpage{61}
$F_{\unsure{o|0}}$ le Module qu'il induit sur $X_o$. On suppose \quad a) Pour tout $s\quad S$, la section $t_s$ induite par $t$ sur la fibre $X_s$ est $\underline O_{X_s}$-régulière et $\underline F_s$-régulière \quad b) $\underline F$ et $\underline O_X$ sont plats par rapport à $S$ \quad c) Pour tout $s\quad S$, $F_{o\ s}$, \deleted{et} $\underline O_{X_{o\ s}}$ \inserted{et $\underline O_{X_s}$} sont de profondeur $\quad2$ en les points fermés de $X_{\deleted{\unsure{o|0}}s\ \deleted{s}}$. Sous ces conditions, il existe un Module cohérent $F$ sur $X$, et un isomorphisme de son complété formel $\widehat F$ le long de $X_o$ avec $\underline F$.

Cet énoncé va résulter du suivant:

\uline{Corollaire} 3.2. \quad Sous les conditions de 3;1. \quad (

\[
\handwritten{\sum_{p\geq0}R^i f_{0*}F_0(q-p)\qquad q-p\geq m}
\]
\[
\handwritten{S_{\unsure{pq|M}}=
\begin{cases}
0&\text{si }p>0\\
\unsure{Sq|S_q}&\text{si }p=0
\end{cases}}
\]
\figure{\begin{tikzpicture}[x=.67cm,y=.67cm,line cap=round,line join=round]
\draw[thin] (-.01,-1.04) .. controls (.08,.62) and (-.16,2.8) .. (-.14,4.31);
\draw[very thin] (-.18,.02)--(-.24,4.10);
\draw[thin] (-.76,.05) .. controls (1.53,.20) and (4.42,.11) .. (6.03,.08);
\draw[very thin] (-.71,-.02)--(-.83,.04)--(-.61,.10);
\draw[thick] (-.02,.05)--(-.04,-.18) .. controls (1.04,-.02) and (3.15,.04) .. (4.67,-.14) .. controls (5.32,-.21) and (5.96,-.30) .. (6.03,-.19) -- (6.01,.18);
\draw[thin] (5.33,-.03)--(5.95,-.09);
\draw[very thin,gray] (.04,.64)--(5.42,.78);
\node[above left] at (-.48,4.04) {$\handwritten{q}$};
\node[right] at (6.04,-.03) {$\handwritten{p}$};
\node at (4.94,3.90) {$\handwritten{F_0(-p)}$};
\begin{scope}[shift={(8.9,-.12)}]
\draw[thin] (0,4.24) .. controls (.22,2.77) and (.04,.79) .. (.10,-1.30);
\draw[thin] (.24,-1.11)--(.09,-2.68);
\draw[thin] (-.9,.10) .. controls (.62,-.06) and (3.1,.15) .. (5.42,.00);
\draw[very thin,gray] (-.95,-3.56)--(6.12,1.89);
\end{scope}
\end{tikzpicture}}
% archive sheet number: 60
